% Nouveau EGA V (cote 25): archive cote 25, site page 165 (PDF page 166).
% typescript. Draft compared: none.
% AI-transcribed by Claude Opus 5.5 (claude-opus-5-5, adaptive thinking, effort high), Message Batches API, batch ega25b (msgbatch_01A7n5YuXubUPNK1zsepi8RG), 2026-11-13; original AI transcription, subsequently checked against both supplied scan parts; see PAGES_CHECKED.tsv and CHECK_DECISIONS.tsv.
\scanpage{25/165}
% NOTE: Both original parts and separated colour-channel contrast views were inspected. This extremely faint forward-reading impression is retained independently; the source transcription of 166 supplied comparison proposals only. Its physical origin is not asserted. All weak readings remain doubtful.
\placedmargin{\unsure{//}} \hand{\unsure{nouveau EGA V}}
\unsure{EGA IV par ?}
\uline{\unsure{Ensemble singulier et supersingulier}.}
% NOTE: Pale curved traces near the left sheet edge and the upper-right spot do not establish an additional textual bracket or a readable folio number.
\unsure{Pages 14 et 15 des notes précédentes, remplacer partout S par Y. Page}
\unsure{15 ligne 10 le} \hand{\unsure{``}}\unsure{i.e.}\hand{\unsure{''}} \unsure{est abusif, c'est la deuxième assertion qui est la}
\unsure{plus précise et celle à laquelle on en a.}
\unsure{La généralisation de prop.7 à laquelle il est fait allusion pages 16}
\unsure{et 17 s'obtient à l'aide du diagramme suivant, qui généralise celui de la}
\unsure{page 4 :}
% NOTE: The horizontal row, upward central arrows and two rising diagonals are visible. The tiny top and bottom zero-like endpoints, several arrowheads and the two diagonal labels remain weak; the proposed readings are retained as doubtful rather than completed from 166.
\diagram{\begin{tikzcd}[ampersand replacement=\&,nodes in empty cells]
 \& \& \unsure{0} \& \& \\
 \& \& \unsure{\hand{\underline{\Omega}}^1_{X/Y}\hand{\otimes}\underline{O}_{V^1}} \arrow[u,handwritten] \& \& \\
\unsure{0} \arrow[r,handwritten] \& \unsure{\underline{N}_{V^1/X}} \arrow[r,handwritten] \arrow[ur,handwritten,"\hand{\unsure{\mu}}"] \& \unsure{\hand{\underline{\Omega}}^1_{X/S}\hand{\otimes}\underline{O}_{V^1}} \arrow[r,handwritten] \arrow[u,handwritten] \& \unsure{\hand{\underline{\Omega}}^1_{V^1/S}} \arrow[r,handwritten] \& \unsure{0} \\
 \& \& \unsure{\hand{\underline{\Omega}}^1_{Y/S}\hand{\otimes}\underline{O}_{V^1}} \arrow[u,handwritten] \arrow[ur,handwritten,"\hand{\unsure{\wr}}"'] \& \& \\
 \& \& \hand{\unsure{0}} \arrow[u,handwritten] \& \&
\end{tikzcd}}
\unsure{où $\underline{N}$ désigne le faisceau conormal. Tous les faisceaux du diagramme sont}
\unsure{localement libres} \unsure{sur $V^1$}\unsure{, et on applique le lemme de la page 7 à $M=\hand{\underline{\Omega}}^1_{X/S}\hand{\otimes}\underline{O}_{V^1}$}
\unsure{et aux deux sous-Modules figurant dans le diagramme précédent. Avec}
\unsure{les notations de ce lemme, on est ramené à prouver que la section a}
 \unsure{de \uline{Hom}(L,P) déduit de} \unsure{$\mu$:} \unsure{$P\hand{\rightarrow}M/Q$ i.e. de} \unsure{$\mu$:} \unsure{$\underline{N}_{V^1/X}\hand{\longrightarrow}\hand{\underline{\Omega}}^1_{X/Y}\hand{\otimes}\underline{O}_{V^1}$,}
\unsure{dont le} \unsure{pré}\unsure{schéma des zéros V(a) est celui qui rend}\hand{\unsure{$\mu$}}\unsure{ strictement de rang}
\hand{\unsure{$<$}}\unsure{d = rang M/Q, satisfait V(a) = V(D), où D est défini au bas de la}
\unsure{page} \unsure{1}\unsure{5. Or les inclusions de SS-préschémas lisses}
\placedmargin{\unsure{$\Vert$}}
\[ \unsure{V^1\hand{\subset}V^o\hand{\subset}X} \]
\unsure{donnent une suite exacte (réf ?)}
\[ \unsure{0\hand{\rightarrow}\underline{N}_{V^o/X}\hand{\otimes}\underline{O}_{V^1}\hand{\rightarrow}\underline{N}_{V^1/X}\hand{\longrightarrow}\underline{N}_{V^1/V^o}\hand{\rightarrow}0\ ,} \]
\unsure{d'autre part on a des isomorphismes canoniques}
\[ \unsure{\underline{N}_{V^o/X}\hand{\simeq}\underline{E}\hand{\otimes}\underline{O}_{V^1}\ ,\quad\underline{N}_{V^1/V^o}\hand{\simeq}\uline{Hom}(W,K)\ ,} \]
\unsure{où W,K sont définis (page 15) comme noyau et conoyau de}
\mspage{\hand{165}}
